% Auto-generated by bin/omega-proof.py
% Source: caostics.txt
\documentclass{jsarticle}
\usepackage[dvipdfmx,final]{graphicx}
\usepackage[utf8]{inputenc}
\usepackage{amsmath,amsthm,amssymb}
\usepackage{geometry}
\geometry{margin=1in}
\begin{document}
   \newcommand{\variint}{%
	\begin{picture}(15,30)
		\put(0,0){
			$ \iint $}
		\put(0,5){\line(15,0){15}}
	\end{picture} %
}
   \newcommand{\vare}{%
	\begin{picture}(5,5)
		\put(4,4){
			 \line(1,1){5} }
		\put(0,0){$\vee$}
	\end{picture} %
}
   \newcommand{\varve}{%
	\begin{picture}(5,5)
		\put(0,0){$\vee$}
		\put(0,5){\line(5,0){5}}
	\end{picture} %
}
\newcommand{\variiint}{%
	\begin{picture}(15,15)
		\put(0,0){
			$ \iiint $}
		\put(0,5){\line(15,0){20}}
	\end{picture} %
}



 \newcommand{\alearletter}{%
	 \begin{picture}(10,10)
		 \put(0,0){
			 $ \square $}
		 \put(0,0){\line(1,1){10}}
	 \end{picture} %
	 }

 \newcommand{\secproduct}{%
	 \begin{picture}(10,20)
		 \put(0,0){
			 $ \nabla $}
		 \put(5,0){+}
	 \end{picture} %
	 }
 \newcommand{\talot}{%
	 \begin{picture}(15,20)
		 \put(0,0){
			 $ \text{\scalebox{2.0}[2]{$\nabla$}}$}
		 \put(5,3){
			 $ \Delta $}
	 \end{picture} %
	 }
 \newcommand{\hazerdnabla}{%
	 \begin{picture}(15,20)
		 \put(0,0){
			 $ \text{\scalebox{2.0}[2]{$\nabla$}}$}
		 \put(5,3){
			 $\text{Y}  $}
	 \end{picture} %
	 }

 \newcommand{\securebox}{%
	 \begin{picture}(15,20)
		 \put(0,0){
			 $ \text{\scalebox{2.0}[2]{$\square$}}$}
		 \put(5,3){
			 $\text{Y}  $}
	 \end{picture} %
	 }


\newcommand{\timebow}{%
	\begin{picture}(5,5)
		\put(0,0){$\text{\rotatebox{90}{$\bowtie$}}$}
	\end{picture} %
	}


 \newcommand{\flowerletter}{%
	 \begin{picture}(10,20)
		 \put(0,0){$\text{\scalebox{1.0}[1]{$\square$}}$}
		 \put(0,0){$\text{\scalebox{1.0}[2]{$\int$}}$}
	 \end{picture} %
	 }
\newcommand{\squareiint}{%
	\begin{picture}(15,30)
		\put(0,0){
			$ \iint $}
		\put(0,5){$\text{\scalebox{2.0}[1]{$\square$}}$}
	\end{picture} %
}

\newcommand{\squareint}{%
	\begin{picture}(15,30)
		\put(0,0){
			$ \int $}
		\put(0,5){$\text{\scalebox{2.0}[1]{$\square$}}$}
	\end{picture} %
}


   \newcommand{\dazanier}{%
	\begin{picture}(12,12)
		\put(0,0){
			$ \bigsqcup $}
		\put(4,3){$\text{\scalebox{1.0}[1]{$\top$}}$}
		\put(4,3){$\text{\scalebox{1.0}[1]{$\bot$}}$}
	\end{picture} %
}
\newcommand{\innabla}{%
	\begin{picture}(13,13)
		\put(0,0){
			$ \nabla $}
		\put(0,0){$\text{\scalebox{2.0}[2]{$\nabla$}}$}
	\end{picture} %
}

\newcommand{\inndelta}{%
	\begin{picture}(20,20)
		\put(0,0){
			$ \Delta $}
		\put(0,0){$\text{\scalebox{2.0}[2]{$\Delta$}}$}
	\end{picture} %
}

\newcommand{\starnabla}{%
	\begin{picture}(20,20)
		\put(3,5){
			$ \Delta $}
		\put(0,0){$\text{\scalebox{2.0}[2]{$\nabla$}}$}
	\end{picture} %
}


  \newcommand{\whitewithbodercircle}
  {\text{\textcircled{\phantom{\textbullet}}}
  \kern-0.9em\text{\textcircled{\phantom{}}}}
     

 
  \newcommand{\whitewithblack}{\text{\textcircled{\textbullet}}}
  \newcommand{\whitewithblacksquare}{\text{\fbox{\text{\rule{0.5em}
  {0.5em}}}}}
  \newcommand{\whitewithwhitesquare}{\text{\fbox{\text{\phantom
  {\rule{0.5em}{0.5em}}}}}}
 \newcommand{\whitewithbodercircleiiint}{%
	 \begin{picture}(20,20)
		 \put(0,0){\scalebox{2.0}[2]{$\whitewithbodercircle$}}
		 \put(0,0){\scalebox{2.0}[2]{$\iiint$}}
	 \end{picture}%
	 }
\newcommand{\whitewithbodercircleint}{%
	 \begin{picture}(20,20)
		 \put(2,2){\scalebox{1.0}[1]{$\whitewithbodercircle$}}
		 \put(2,2){\scalebox{1.0}[2]{$\int$}}
	 \end{picture}%
	 }
\newcommand{\whitewithbodercircleiint}{%
	 \begin{picture}(20,20)
		 \put(0,0){\scalebox{1.0}[1]{$\whitewithbodercircle$}}
		 \put(2,2){\scalebox{1.0}[2]{$\iint$}}
	 \end{picture}%
	 }
 
 \newcommand{\whitewithbodercirclevarint}{%
	 \begin{picture}(13,13)
		 \put(0,0){\scalebox{2.0}[2]{$\whitewithbodercircle$}}
		 \put(0,0){\scalebox{2.0}[2]{$\varint$}}
	 \end{picture}%
	 }
 
 \newcommand{\whitewithbodercirclevariiint}{%
	 \begin{picture}(13,13)
		 \put(0,0){\scalebox{2.0}[2]{$\whitewithbodercircle$}}
		 \put(0,0){\scalebox{2.0}[2]{$\variiint$}}
	 \end{picture}%
	 }
 
 \newcommand{\average}{%
	  \begin{picture}(30,45)
		  \put(0,0){\scalebox{2.0}[2]{$\square$}}
		  \put(4,4){$\setminus$}
	  \end{picture}%
  } 
  \newcommand{\smallaverage}{%
	  \begin{picture}(30,45)
		  \put(0,0){\scalebox{1.0}[1]{$\square$}}
		  \put(0,0){$\setminus$}
	  \end{picture}%
  } 
 
  \newcommand{\whitewithboder}{%
	  \begin{picture}(15,30)
		  \put(0,0){\scalebox{2.0}[2]{$\square$}}
		  \put(4,4){$\square$}
	  \end{picture}%
  } 
  \newcommand{\whitewithnabla}{%
	  \begin{picture}(15,30)
		  \put(0,0){\scalebox{2.0}[2]{$\nabla$}}
		  \put(3,3){$\nabla$}
	  \end{picture}%
  } 



\section*{主題 / Topic}
$$ \Phi(\delta(H\Psi))^{\nabla} = \scalebox{1}[2]{\variiint} \mathrm{cohom D\chi}(M) $$。

$$ \varve |: \chi \to H\Psi^{D} \to \bigoplus {i\hbar} \otimes x_1 $$

\section*{定理 / Theorem}
$$ \Phi(\delta(H\Psi))^{\nabla} = \scalebox{1}[2]{\variiint} \mathrm{cohom D\chi}(M) $$。

$$ \varve |: \chi \to H\Psi^{D} \to \bigoplus {i\hbar} \otimes x_1 $$

Under equations are circle element of step with neipa number represent。

and circle element of pai also zeta function from logment system.

This number system express imaginary of reverse pole in circle function。

\section*{予想 / Conjecture}
$$ \Phi(\delta(H\Psi))^{\nabla} = \scalebox{1}[2]{\variiint} \mathrm{cohom D\chi}(M) $$

$$ \varve |: \chi \to H\Psi^{D} \to \bigoplus {i\hbar} \otimes x_1 $$

\section*{証明 / Proof}
$$ \Phi(\delta(H\Psi))^{\nabla} = \scalebox{1}[2]{\variiint} \mathrm{cohom D\chi}(M) $$

$$ \varve |: \chi \to H\Psi^{D} \to \bigoplus {i\hbar} \otimes x_1 $$

Under equations are circle element of step with neipa number represent。

and circle element of pai also zeta function from logment system.

This number system express imaginary of reverse pole in circle function。

After all, this system of neipa and pai number of step function。

\section*{結論 / Conclusion}
this express with function escourt into rout of g conclude with。

These equation are concluded of being formula,。

$$ \Phi(\delta(H\Psi))^{\nabla} = \scalebox{1}[2]{\variiint} \mathrm{cohom D\chi}(M) $$

\end{document}
